%==============================================================================
%  BAB IV - KESIMPULAN DAN SARAN
%==============================================================================
\section{Kesimpulan dan Saran}

\subsection{Kesimpulan}

Milestone 4 berhasil menyatukan seluruh komponen kompiler Arion yang dibangun pada
milestone sebelumnya menjadi sebuah kompiler yang utuh dan dapat dieksekusi. Dari
pengerjaan tahap ini dapat ditarik beberapa kesimpulan.

Pertama, \emph{Decorated AST} hasil analisis semantik berhasil diratakan menjadi
\emph{Three-Address Code} melalui penelusuran berbasis DFS. Ekspresi diterjemahkan
secara postorder sehingga operand selalu tersedia di stack sebelum operatornya
dijalankan, sementara struktur control flow \texttt{if-else} dan \texttt{while}
diratakan menggunakan kombinasi label dan instruksi lompat dengan konvensi
\emph{IF\_FALSE}. Intermediate Code yang dihasilkan terbukti identik dengan contoh
pada spesifikasi.

Kedua, interpreter berbasis \emph{stack machine} mampu mengeksekusi Intermediate
Code dengan benar melalui siklus \emph{fetch--decode--execute}, mengelola memori
dalam bentuk stack frame (static link, dynamic link, return address, dan variabel),
serta menghasilkan keluaran nyata program melalui operasi \texttt{WRT}/\texttt{WRTLN}.

Ketiga, aspek keamanan runtime---termasuk tantangan bonus---telah
diimplementasikan. Interpreter menangani pembagian dengan nol, overflow/underflow
aritmatika, stack overflow/underflow, stack corruption, stack smashing, akses
memori dan target lompat di luar batas, serta pemeriksaan tipe saat eksekusi.
Seluruh kondisi error dilaporkan secara \emph{graceful} dengan pesan informatif
tanpa membuat program \emph{crash}, sebagaimana ditunjukkan pada Kasus 5 dan 6.

Keempat, kode disusun secara modular dengan pemisahan tanggung jawab yang jelas
antara pembentukan kode (\texttt{intermediate/}) dan eksekusi
(\texttt{interpreter/}), dijembatani oleh pola \emph{adapter} sehingga kedua sisi
tetap terpisah (\emph{loosely coupled}).

\subsection{Saran}

Beberapa hal dapat dikembangkan lebih lanjut. Implementasi pemanggilan
fungsi/prosedur (\texttt{CAL}/\texttt{RET} dengan \emph{nested scope} dan
\emph{passing} argumen penuh) masih bersifat dasar dan dapat dilengkapi agar
mendukung rekursi serta variabel lokal antar-frame secara menyeluruh. Selain itu,
pembentukan kode dapat diperluas untuk menangani tipe data terstruktur seperti
\emph{array} dan \emph{record} beserta validasi indeks \emph{out-of-bounds} saat
runtime. Dari sisi optimasi, \emph{constant folding} dan eliminasi instruksi
redundan dapat ditambahkan pada tahap pembentukan TAC untuk menghasilkan kode yang
lebih ringkas. Terakhir, cakupan pengujian otomatis dapat diperluas dengan kerangka
uji regresi agar setiap perubahan dapat diverifikasi dengan cepat.
